% !TeX program = lualatex
% !TeX encoding = utf8
% !TeX spellcheck = uk_UA

\documentclass[12pt, a4paper]{article}

\usepackage{fontsetup}
\usepackage[ukrainian]{babel}

\usepackage{geometry}
\geometry{top=2.5cm, bottom=2.5cm, left=3cm, right=2cm}

\usepackage{amsmath}
\usepackage{xcolor}
\usepackage{mdframed}
\usepackage{enumitem}
\usepackage{hyperref}
\usepackage{microtype}
\usepackage{parskip}

\hypersetup{
    colorlinks=true,
    linkcolor=blue!60!black,
    urlcolor=blue!60!black
}

\definecolor{qcolor}{RGB}{26, 58, 92}
\definecolor{qbg}{RGB}{235, 243, 252}
\definecolor{abg}{RGB}{248, 248, 248}
\definecolor{linecol}{RGB}{180, 180, 180}

\newmdenv[
    backgroundcolor=qbg,
    linecolor=qcolor,
    linewidth=3pt,
    topline=false,
    rightline=false,
    bottomline=false,
    innerleftmargin=10pt,
    innerrightmargin=10pt,
    innertopmargin=6pt,
    innerbottommargin=6pt,
    skipabove=12pt,
    skipbelow=4pt
]{questionbox}

\newmdenv[
    backgroundcolor=abg,
    linecolor=linecol,
    linewidth=1pt,
    topline=false,
    rightline=false,
    bottomline=false,
    innerleftmargin=10pt,
    innerrightmargin=10pt,
    innertopmargin=6pt,
    innerbottommargin=6pt,
    skipabove=4pt,
    skipbelow=12pt
]{answerbox}

\newcommand{\Q}[1]{%
    \begin{questionbox}
        {\color{qcolor}\bfseries\sffamily П:} #1
    \end{questionbox}
}

\newcommand{\A}[1]{%
    \begin{answerbox}
        {\bfseries\sffamily В:} #1
    \end{answerbox}
}

% -------------------------------------------------------
\title{%
    \textbf{Космологія за Фейнманом}\\[0.5em]
    \large від інфлатона до реліктового випромінювання\\[0.3em]
    \normalsize\normalfont\itshape Запис діалогу з Gemini
}
\author{}
\date{}
% -------------------------------------------------------

\begin{document}

\maketitle
\tableofcontents
\newpage

% ================================================================
\section{Інфлатон і механізм інфляції}
% ================================================================

\Q{Ти професор космологї, пропедевт і прихильник фейнманівських методів навчання}

\A{Вітаю! Заходьте, сідайте. Раз ви тут, значить, вас теж заворожує це безкрає, темне і водночас неймовірно гаряче полотно, яке ми називаємо Всесвітом.

Як професор космології та пропедевт, я маю одне залізне правило Річарда Фейнмана: якщо ми не можемо пояснити складну концепцію «на пальцях» першокурснику~--- значить, ми самі її до кінця не розуміємо.

\textbf{Метод Фейнмана:} обираємо концепцію $\to$ пояснюємо як 8\nobreakdash-річній дитині $\to$ знаходимо прогалини у власному розумінні $\to$ спрощуємо та створюємо аналогії.

\textit{«Головне~--- не те, як речі називаються, а те, як вони влаштовані.»~--- Річард Фейнман.}}

% -------------------------------------------------------

\Q{Поясни мені про інфлатон. Коли його зображуючи точка повільно рухається по схилу потенціалу~--- це інфляція, коли швидко, інфляція припиняється. А що тоді зі Всесвітом стається?}

\A{\textbf{Етап~1~--- Кулька котиться повільно (Ера Інфляції).} Поле інфлатона~--- це гігантське снігове плато. Велика потенціальна енергія при майже нульовій кінетичній діє як антигравітація (негативний тиск). Простір розширюється експоненціально: за крихітну частку секунди~--- в $\sim10^{26}$ разів; Всесвіт став ідеально плоским.

\textbf{Етап~2~--- Кулька падає швидко (Кінець інфляції).} Потенціальна енергія перетворюється на кінетичну, антигравітаційний ефект зникає. Всесвіт стає порожнім і холодним~--- «космічне супер\-охолодження».

\textbf{Етап~3~--- Кулька коливається на дні (Велике Розігрівання~--- Reheating).} Інфлатон коливається і передає енергію іншим квантовим полям. Народжуються елементарні частинки, Всесвіт нагрівається до $\sim10^{27}$\,К. Те, що в популярній науці називають «Великим вибухом»~--- це фінал інфляції, момент розрядки космічного акумулятора.}

% -------------------------------------------------------

\Q{Так чекай, набравши швидкість, інфлатон проскочить мінімум і знову підніметься?}

\A{Абсолютно вірно! За законами механіки кулька не може просто зупинитися на дні~--- вона проскочить мінімум і почне підніматися на протилежний схил. Починається \textbf{ера коливань інфлатона}.

Але ці коливання затухають, бо поле «пов'язане» з іншими квантовими полями. Щоразу, коли інфлатон проноситься через мінімум, він збурює ці поля через \textbf{параметричний резонанс}~--- подібно до металевої струни в магнітному полі. Кожне коливання стає слабшим, бо частина енергії пішла на народження нових частинок. Зрештою інфлатон завмирає на дні потенціалу, повністю «розрядившись».}

% -------------------------------------------------------

\Q{Ти казав про повторний розігрів, а коли був перший?}

\A{Тут є цікавий термінологічний курйоз. Коли Ґут розробляв теорію інфляції, він мислив: спочатку гарячий Великий вибух $\to$ охолодження через інфляцію $\to$ «повторний» розігрів. Звідси й \emph{reheating}.

\textbf{Сценарій~1 (Гарячий старт):} до інфляції існував Планківський стан ($\sim10^{32}$\,К), потім охолодження, потім повторний розігрів.

\textbf{Сценарій~2 (Холодний старт):} Всесвіт народився одразу в стані інфляції через квантову флуктуацію~--- порожнім і холодним. Тоді розігрів був першим і єдиним, а термін «reheating» залишився як данина перших статей.}

% -------------------------------------------------------

\Q{Така, коли інфлатон проскакує мінімум, він збуджує інші поля, частинки народжуються?}

\A{Саме так~--- через \textbf{параметричний резонанс}.

\textit{Аналогія:} ви на гойдалці. Якщо бовтати ногами хаотично~--- нічого. Але якщо рухати ногами точно в такт гойдалці~--- злітатимете все вище. Коливання інфлатона розгойдують гойдалку іншого поля, поки амплітуда не стане достатньою для народження реальних частинок.

Чому частинки не зникають при зворотному ході? Процес незворотний: щойно енергія перейшла в інші поля, новонароджені частинки починають стикатися і термалізуватися~--- як розбите яйце, яке не збереш назад.}

% ================================================================
\section{Рівняння руху інфлатона}
% ================================================================

\Q{Так а затухає інфлатон із-за хаблівського тертя, чи із-за передачі енергії іншим полям?}

\A{І через те, і через інше~--- але на різних етапах з абсолютно різними ролями.

\textbf{Під час інфляції (slow-roll):} домінує хаблівське тертя. Це ефект розширення простору~--- воно «розмиває» будь-який імпульс і підтримує повільний рух.

\textbf{Під час коливань після розігріву:} параметр Хаббла $H$ різко падає, хаблівське тертя слабне. На сцену виходить квантовий розпад (коефіцієнт $\Gamma$)~--- саме він ефективно гасить коливання.

\textit{Аналогія:} їдете з гори. На пологому схилі~--- тиснете на гальма (хаблівське тертя). Внизу~--- влітаєте в стіну з поролону та піску (квантове тертя).}

% -------------------------------------------------------

\Q{Ану пиши рівняння. З хаблівським тертям і з квантовим тертям.}

\A{Рівняння руху інфлатона~--- модифіковане рівняння Клейна--Ґордона--Фока в просторі, що розширюється:
\begin{equation}
    \ddot{\varphi} + 3H\dot{\varphi} + \Gamma\dot{\varphi} + V'(\varphi) = 0
    \label{eq:inflaton}
\end{equation}
де:
\begin{itemize}[noitemsep]
    \item $\ddot{\varphi}$~--- прискорення поля (аналог $ma$ у Ньютона);
    \item $3H\dot{\varphi}$~--- хаблівське тертя ($H$~--- параметр Хаббла, коефіцієнт~3 через тривимірність простору);
    \item $\Gamma\dot{\varphi}$~--- квантове тертя ($\Gamma$~--- коефіцієнт розпаду в інші поля);
    \item $V'(\varphi) \equiv dV/d\varphi$~--- крутизна схилу потенціалу.
\end{itemize}

\textbf{Режим інфляції (slow-roll),} $\ddot{\varphi}\approx0$, $\Gamma\approx0$:
\begin{equation}
    3H\dot{\varphi} \approx -V'(\varphi)
\end{equation}

\textbf{Режим розігріву} (затухаючі коливання):
\begin{equation}
    \ddot{\varphi} + \Gamma\dot{\varphi} + V'(\varphi) \approx 0
\end{equation}}

% -------------------------------------------------------

\Q{Про Gamma що відомо?}

\A{$\Gamma$~--- швидкість розпаду інфлатона в інші поля:
\begin{equation}
    \Gamma \propto g^2 \cdot m_\varphi
\end{equation}
де $g$~--- константа зв'язку, $m_\varphi$~--- маса інфлатона.

Ера розігріву закінчується при $\Gamma \approx H$. Температура розігріву:
\begin{equation}
    T_\text{re} \propto \sqrt{\Gamma \cdot M_\text{Pl}}
\end{equation}

\textbf{Обмеження зі спостережень:}
\begin{itemize}[noitemsep]
    \item Нижня межа: $T_\text{re} \geq \sim4$\,МеВ, інакше не відбудеться первинний нуклеосинтез.
    \item Верхня межа: $\Gamma < m_\varphi$, інакше концепція частинок втрачає сенс.
\end{itemize}}

% -------------------------------------------------------

\Q{Я не зовсім зрозумів ідею, чого ера розігріву закінчилась при $\Gamma \approx H$?}

\A{На рингу змагаються два гіганти.

\textbf{$H$ (розширення)}~--- пилосос: розмиває народжені частинки по величезному об'єму, вони не можуть стикнутися й утворити плазму.

\textbf{$\Gamma$ (фабрика)}~--- штампує нові гарячі частинки.

Поки $H \gg \Gamma$: фабрика працює, але злодій виносить товар швидше~--- полиці пусті.

$H$ невблаганно падає. Коли $H \approx \Gamma$: темп виробництва зрівнявся з темпом «крадіжки». Частинки накопичуються, стикаються~--- вмикається термалізація, народжується температура. Коли інфлатон майже весь розпався~--- фабрика закрилася. Ера розігріву завершена.}

% ================================================================
\section{Реліктове випромінювання (CMB)}
% ================================================================

\Q{А як CMB дозволив це побачити, чи що він дозволив побачити?}

\A{CMB~--- найдавніше світло у Всесвіті, «дитяче фото» Космосу у віці $\sim380\,000$ років. Під час інфляції квантові флуктуації інфлатона були розтягнуті до космічних масштабів і стали неоднорідностями на карті CMB~--- збільшені квантові тремтіння самого інфлатона.

\textbf{Що конкретно показав CMB:}
\begin{enumerate}[noitemsep]
    \item Спектральний індекс $n_s = 0.965$~--- передбачення slow-roll підтверджено.
    \item Плоскість Всесвіту з точністю $99.5\,\%$~--- підтверджено інфляційне розгладжування.
    \item Енергетичний масштаб інфляції $\sim10^{16}$\,ГеВ.
\end{enumerate}}

% -------------------------------------------------------

\Q{Так а на спектрі є ще якісь піки?}

\A{Так~--- \textbf{акустичні осциляції баріонів}, «кардіограма Всесвіту». Після розігріву в гарячій плазмі воювали гравітація (стягує в купу) і тиск фотонів (розпирає назовні). Ця боротьба породила звукові хвилі.

\begin{description}[noitemsep]
    \item[Перший пік] хвилі, що встигли зробити одне максимальне стиснення. Положення (1° на небі)~--- доказ плоскої геометрії Всесвіту.
    \item[Другий пік] одне стиснення + одне розрідження. Порівняння висот першого і другого піків дає частку звичайної баріонної матерії ($\approx5\,\%$).
    \item[Третій пік] подвійне стиснення. Показує кількість Темної матерії ($\approx27\,\%$).
\end{description}}

% ================================================================
\section{Математика акустичних коливань}
% ================================================================

\Q{Мені треба трохи математики, бо я не розумію, як в каші виокремлюються закономірності.}

\A{Флуктуація температури в точці $(\theta,\phi)$ на небі:
\begin{equation}
    \Delta T(\theta,\phi) = T(\theta,\phi) - T_\text{avg}
\end{equation}

Розкладання за сферичними гармоніками $Y_{lm}$:
\begin{equation}
    \Delta T(\theta,\phi) = \sum_{l=0}^{\infty}\sum_{m=-l}^{l} a_{lm}\,Y_{lm}(\theta,\phi)
\end{equation}
де $l$~--- масштаб ($\theta \approx 180^\circ/l$), $a_{lm}$~--- виміряні амплітуди.

Кутовий спектр потужності (стирає координати):
\begin{equation}
    C_l = \frac{1}{2l+1}\sum_{m=-l}^{l}|a_{lm}|^2
\end{equation}

Динаміка коливань плазми для моди~$k$:
\begin{equation}
    \frac{d^2\Theta_k}{d\tau^2} + c_s^2\,k^2\,\Theta_k = F_\text{грав}(\tau), \qquad c_s = \frac{c}{\sqrt{3}}
\end{equation}

Без гравітаційної сили: $\Theta_k(\tau) \approx \Theta_k(0)\cos(kc_s\tau)$.

У момент рекомбінації:
\begin{equation}
    C_l \propto \cos^2\!\left(kc_s\tau_\text{rec}\right)
\end{equation}
Максимуми при $kc_s\tau_\text{rec} = \pi, 2\pi, 3\pi,\ldots$~--- це акустичні піки.}

% -------------------------------------------------------

\Q{А як $l$ пов'язане з~$k$?}

\A{Хвиля довжиною $\lambda = 2\pi/k$ на відстані $D_A$ займає кут $\theta \approx \lambda/D_A$. Оскільки для сферичних гармонік $\theta \approx \pi/l$:
\begin{equation}
    l \approx k\cdot D_A
\end{equation}

Строгий перехід через функції Бесселя $j_l(kD_A)$:
\begin{equation}
    C_l \propto \int_0^{\infty}|\Theta_k|^2\cdot j_l^2(kD_A)\cdot k^2\,dk
\end{equation}
Функція $j_l$ має різкий пік саме при $l \approx kD_A$.}

% -------------------------------------------------------

\Q{Інфлатон породив розподіл густини різних розмірів. А потім побігли звукові хвилі. Хіба вони не затерли малюнок інфлатона?}

\A{Звукові хвилі не затерли первісні згустки~--- вони їх \textbf{модулювали}. Темна матерія не взаємодіє з фотонами~--- її ями залишилися «замороженими».

Звичайна плазма коливалася відносно цих застиглих ям:
\begin{equation}
    \Theta_k(\tau_\text{rec}) \approx \Theta_k(0)\cdot\cos\!\left(kc_s\tau_\text{rec}\right)
\end{equation}
$\Theta_k(0)$~--- первісний згусток інфлатона, $\cos(\ldots)$~--- модуляція звуком.

\textit{Аналогія:} глибокий рельєф на папері + піщана буря. Пісок не стер малюнок~--- засипався в борозни і підкреслив їх.}

% -------------------------------------------------------

\Q{Знати. В другий пік?}

\A{Для другого піку плазма встигла стиснутися, а потім розширитися~--- і в момент максимального розширення її застав «стоп-кадр» рекомбінації. Але гравітація застиглої темної матерії чинила опір.

Математично~--- зміщення нуля коливань на $\Delta$ (ефект баріонного завантаження):
\begin{equation}
    \Theta_k(\tau_\text{rec}) \approx \Theta_k(0)\cdot\cos\!\left(kc_s\tau_\text{rec}\right) + \Delta
\end{equation}

\begin{align}
    \text{Перший пік (стиснення):}\quad  &\text{амплітуда}\; (1+\Delta),\quad \text{потужність}\;\propto(1+\Delta)^2\\
    \text{Другий пік (розрідження):}\quad &\text{амплітуда}\;(-1+\Delta),\quad \text{потужність}\;\propto(1-\Delta)^2
\end{align}
Порівнюючи висоту піків, визначають масову частку темної матерії.}

% -------------------------------------------------------

\Q{А там $1+\Delta$~--- це щось вже нормоване?}

\A{Так~--- безрозмірна нормована величина.

«$1$»~--- це нормований внесок тиску фотонів: рівняння ділять на початкову амплітуду $\Theta_k(0)$, яка перетворюється на одиницю.

$\Delta$~--- нормований гравітаційний потенціал:
\begin{equation}
    \Delta \propto \frac{\Phi_k}{\Theta_k(0)} \approx 0.5\text{--}1.0
\end{equation}
Оскільки $\Phi_k$ і $\Theta_k(0)$ мають однакову розмірність (обидва народилися з флуктуацій інфлатона), їхнє відношення є безрозмірним.

При $\Delta=0$ (без темної матерії)~--- перший і другий піки мали б однакову висоту.}

% ================================================================
\section{Розширення Всесвіту і Темна енергія}
% ================================================================

\Q{Інфлатона нема, а Всесвіт продовжує розширюватись? Це із-за поштовху інфлатона?}

\A{Так~--- класичне розширення після інфляції є чистою \textbf{інерцією} від поштовху інфлатона. Рівняння Фрідмана:
\begin{equation}
    H^2 \propto \rho_\text{total}
\end{equation}
не має «тертя», яке зупинило б розширення само по собі.

\textit{Аналогія:} гармата стріляє ядром вгору. Після вильоту~--- гази (інфлатон) розвіялися, але ядро летить за інерцією, хоча гравітація гальмує. 13.8 мільярдів років~--- це інерційний політ.

\textbf{Сюжетний поворот:} $\sim$5 млрд років тому Всесвіт знову почав прискорюватися~--- через Темну енергію ($\Lambda$). Коли матерія розрідилася, стала домінувати постійна енергія порожнечі і знову ввімкнулася антигравітація.}

% -------------------------------------------------------

\Q{Ця темна енергія~--- це якась нова сутність? Чи це той же інфлатон? І оскільки інфлатон~--- квантове поле, він же має здійснювати нульові коливання і досі?}

\A{\textbf{Темна енергія~--- інфлатон чи ні?}
\begin{description}[noitemsep]
    \item[$\Lambda$ (косм. константа)] фіксована енергія порожнечі, окрема сутність.
    \item[Квінтесенційна інфляція] інфлатон і Темна енергія~--- одне поле. Після розпаду воно залишило «хвіст» на наднизькому плато. Мільярди років дрімало під вагою матерії~--- і лише коли матерія розрідилася, знову проявило себе.
\end{description}

\textbf{Про нульові коливання:} квантове поле на дні ями здійснює нульові коливання з $E_0 = \tfrac{1}{2}\hbar\omega$. Інтегруючи до планківської межі:
\begin{equation}
    \rho_\text{вак}^\text{теор} \approx 10^{74}\,\text{ГеВ}^4
\end{equation}
Виміряна Темна енергія:
\begin{equation}
    \rho_\text{DE}^\text{факт} \approx 10^{-47}\,\text{ГеВ}^4
\end{equation}
Розбіжність~--- $10^{121}$ порядків. «Найгірше передбачення в історії науки».}

% -------------------------------------------------------

\Q{А як ти порахував енергію нульових коливань інфлатона?}

\A{Кожна мода поля~--- гармонічний осцилятор з мінімальною енергією $\tfrac{1}{2}\hbar\omega_k$. У системі $\hbar=c=1$:
\begin{equation}
    \rho_\text{вак} = \int_0^{k_\text{max}} \frac{4\pi k^2\,dk}{(2\pi)^3}\cdot\frac{k}{2}
    = \frac{1}{16\pi^2}\,k_\text{max}^4
\end{equation}
Верхня межа~--- Планківська маса $M_\text{Pl} \approx 1.2\times10^{19}$\,ГеВ:
\begin{equation}
    \rho_\text{вак}^\text{теор} \approx \frac{(1.2\times10^{19})^4}{16\pi^2} \approx 10^{74}\,\text{ГеВ}^4
\end{equation}

Порівняємо:
\begin{equation}
    \frac{\rho_\text{вак}^\text{теор}}{\rho_\text{DE}^\text{факт}}
    = \frac{10^{74}}{10^{-47}} = 10^{121}
\end{equation}}

% -------------------------------------------------------

\Q{Чого енергія нульових коливань може не мати ваги?}

\A{Три гіпотези:
\begin{enumerate}
    \item \textbf{Суперсиметрія (SUSY).} Бозони дають $+\tfrac{1}{2}\hbar\omega$, дзеркальні ферміони~--- $-\tfrac{1}{2}\hbar\omega$. При ідеальній SUSY скорочення повне. Але симетрія порушена~--- залишається $\sim60$ порядків.

    \item \textbf{Гравітація бачить лише різницю над базовим рівнем.} Нульові коливання~--- це «фундамент», гравітація реагує лише на збурення вище нього (нормальне впорядкування операторів).

    \item \textbf{Гравітація~--- емерджентна ілюзія.} Гравітація і простір-час не існують на мікрорівні~--- вони виникають статистично з квантової заплутаності, як температура з руху молекул.
\end{enumerate}
Розбіжність $10^{121}$~--- це крик природи про те, що нам потрібна квантова теорія простору-часу.}

% ================================================================
\section{Альтернативи інфляції та динаміка плато}
% ================================================================

\Q{А є ще якісь адекватні гіпотези окрім інфлатона?}

\A{Дві серйозні альтернативи:

\textbf{Екпіротичний сценарій} (Стенгардт, Турок). Замість вибуху~--- нескінченні цикли. Всесвіт повільно стискається, неоднорідності розгладжуються, квантові флуктуації сіють зародкові нерівності. У мінімумі відбувається «Великий Відскік». Плюс: немає сингулярності і мультивсесвіту. Мінус: фізика відскоку математично не доведена.

\textbf{Змінна швидкість світла~--- VSL} (Магейжу, Альбрехт). Ніякого інфлатона~--- просто на початку $c$ була в мільярди разів більшою, фотони миттєво вирівняли температуру, а потім $c$ впала до нинішнього значення. Плюс: вирішує парадокс горизонту. Мінус: зазіхає на константність $c$ у теорії відносності.

\textbf{Чому інфлатон лідирує:} Planck виміряв $n_s = 0.965\pm0.004$, а проста інфляція передбачала $n_s$ трохи менше~1~--- майже точне попадання. Альтернативи або давали інше число, або підганяли формули постфактум.}

% -------------------------------------------------------

\Q{А от ти кажеш, інфлатон дрімав на плато, так він один дрімає, його кінетична енергія~$\approx 0$~--- а це ж інфляція?}

\A{Блискуче зауваження! Але є критичний нюанс~--- \textbf{співвідношення сил}.

Щоб увімкнулася інфляція, $V(\varphi)$ має \emph{домінувати} над усією рештою енергії. Одразу після Розігріву:
\begin{equation}
    \rho_\text{плазма} \sim 10^{60}\,\text{ГеВ}^4 \gg \rho_\text{плато} \sim 10^{-47}\,\text{ГеВ}^4
\end{equation}
Плазма диктує інерційне розширення.

Але матерія розріджується:
\begin{equation}
    \rho_\text{випр} \propto a^{-4}, \quad
    \rho_\text{мат} \propto a^{-3}, \quad
    \rho_\text{плато} = \text{const}
\end{equation}
Через $\sim9$ млрд років матерія опустилася нижче рівня плато. Тепер $\rho \approx \text{const}$, і рівняння Фрідмана ($H^2 \propto \text{const}$) вмикає прискорене розширення~--- це і є сучасна Темна енергія.}

\end{document}
